\documentclass[a4paper,10pt]{article}
\usepackage[francais]{babel}
\usepackage{url}

\input{letter-ensiie.sty}

\enCouleur

\begin{document}

\name{St\'ephane \textsc{Genaud}}
\fonction{Professeur ENSIIE, mise à disposition Université de Strasbourg}
\telephoneillkirch{+33 (0)3688 54542}
\email{genaud@ensiie.fr}

\begin{letter}[Attestation et Recommendation de Pierre \textsc{Boutry}]{à}%
{CNU 27, Campagne de Qualification 2020}

Chers Collègues,

En tant que directeur de l'antenne de l'ENSIIE à Strasbourg de 2013 à 2107,
j'atteste par ce document que Pierre Boutry a assuré les enseignements suivants
auprès des élèves-ingénieurs de ce cursus.

\noindent 
Quantitavement, Pierre Boutry a enseigné, dans le cadre de son contrat doctoral, environ 64h (toutes les heures
sont en heures équivalent TD):
\begin{itemize}
	  \item  en 2013-2014: Analyse Numérique (35h), Recherche Opérationnelle (28h)
	  \item  en 2014-2015: Analyse Numérique (56h), Logique (10,5h)
	  \item  en 2015-2016:  Analyse Numérique (21h), Logique (21h), Programmation impérative (28h)
\end{itemize}

\smallskip

\noindent Il a ensuite occupé un poste d'ATER en effectuant un service complet. Il a assuré les
enseignements suivants ainsi que des encadrements de projets étudiants.
\begin{itemize}
	  \item  en 2016-2017:  Programmation Fonctionnelle (79h), Logique (21h), 
		    Outils Mathématique (52,5h), Ingénierie de la Preuve (21h)
\end{itemize}

\bigskip

\noindent Qualitativement, je tiens à souligner le très grand professionnalisme et l'investissement total
de Pierre Boutry dans ses enseignements. 
%
En 2013-2014, 2014-2015 et 2015-2016, il a par exemple accepté de prendre la charge du 
cours d'analyse numérique. Je peux également témoigner qu'il s'est investi à de nombreuses
reprises pour organiser des séances informelles pour ré-expliquer aux étudiants volontaires
ce qui n'avait pas été compris lors de ses cours. Humainement, Pierre Boutry est 
extrêment dévoué et consciencieux.

Je veux également souligner la polyvalence technique de Pierre Boutry.
Comme le montre la liste des enseignements effectués, il est capable d'enseigner l'algorithmique
et la programmation (les langages support étaient C et Caml) comme des matières aux
fondements mathématiques (Logique, Recherche Opérationnelle, Preuve). 

En conclusion, Pierre Boutry a été un élément extrêmement précieux pour nos enseignements,
et je le recommande au plus haut point sur le plan de l'enseignement.

Je vous prie, chers et chères collègues, d'agréer mes sentiments respectueux.
\end{letter}
\begin{flushright}
\includegraphics[width=.2\textwidth]{signgenaud.jpg}
\end{flushright}
\end{document}
